% !TeX root = ../main.tex
答辩委员会听取了答辩人彭森就学位论文《周期级多核神经网络芯片模拟器》所做的报告与答辩，经质询与讨论，认为论文针对多核神经网络芯片在设计阶段对端到端周期级评估与瓶颈归因的需求，给出了以调度器生成中间表示为输入、在统一时间轴上将计算核心、片上网络与片外存储耦合推进的模拟器总体方案。阐述了统一任务抽象、接口化松耦合的子模块后端集成方式以及多时钟域同步等关键机制。并通过微基准与硬件模拟对照、真实负载上的阶段分解与设计空间探索等实验，验证了模拟器的有效性与实用价值。论文选题紧密结合集成电路工程领域前沿与应用需求，工作饱满，论述清楚，层次分明。

论文取得的主要创新性成果包括：

\begin{enumerate}[wide]
    \item 构建了面向多核神经网络芯片的端到端周期级模拟框架，实现了从解析、周期推进到统计与踪迹输出的完整链路。
    \item 提出了统一任务表示与接口化松耦合集成结构，为论文复现和进一步研究开发提供了基础。
    \item 建立了较为系统的实验与分析方法，既证明了工作的可用性，又提供了新的瓶颈分析和软硬件协同优化方法。
\end{enumerate}

论文工作表明答辩人在集成电路工程方面具有较扎实的基础理论和专业知识，具备独立从事科研与工程实现的能力，论文写作规范，图表与引用较为完整，答辩过程中能够较好地回答委员会提出的问题，达到硕士学位论文要求。

答辩委员会表决，一致同意通过论文答辩，并建议授予答辩人工程硕士学位。

% 提交前请将上文中的「答辩人」替换为本人姓名；若学院决议文本对学位表述有固定模板，请以学院最新版为准微调末句措辞。
